% ============================
% Pre-load options
% ----------------------------
% xcolor must receive [table] BEFORE any package (tcolorbox, tikz, ...)
% pulls it in without options. Otherwise, "\usepackage[table]{xcolor}"
% later is a silent no-op under pdflatex but triggers an option-clash
% cascade under tex4ht ("Paragraph ended before \@reset@ptions").
% ============================
\PassOptionsToPackage{table}{xcolor}

% ============================
% Engine / output mode detection
% ----------------------------
% \ifPDFTeX  - true under pdflatex (also under tex4ht's htlatex driver)
% \ifLuaTeX  - true under lualatex
% \ifXeTeX   - true under xelatex
% \ifSEhtml  - true when being processed by tex4ht / make4ht
%             (named "SEhtml" to avoid clashing with package internals
%              in the @-letter namespace; the @ form was being reset
%              somewhere downstream, silently flipping the flag back.)
% ============================
\usepackage{iftex}

\newif\ifSEhtml
\SEhtmlfalse
\makeatletter
\@ifundefined{HCode}{}{\SEhtmltrue}
\makeatother

% ============================
% Geometry
% ============================
\usepackage[
  letterpaper,
  left=0.5in,
  right=0.5in,
  top=0.5in,
  bottom=1.0in
]{geometry}

% ============================
% Required packages (always)
% ============================
\usepackage[most]{tcolorbox}
\usepackage{booktabs}
\usepackage{fancyhdr}
\usepackage{graphicx}
\usepackage{lastpage}
\usepackage{listings}
\usepackage{setspace}
\usepackage{subcaption}
\usepackage{tabularx}
\usepackage{tabularray}
\usepackage{tikz}
\usepackage{xcolor}             % [table] is set via PassOptionsToPackage above
\usepackage{xparse}
\usepackage[siunitx, american, RPvoltages]{circuitikz}
\usepackage{array}
\usepackage{colortbl}
\usepackage{pdflscape}
\usepackage{pgfplots}
\usepackage{xurl}
\usetikzlibrary{arrows.meta, angles, quotes, shapes.symbols, arrows.meta}

% ============================
% minted
% ----------------------------
% In HTML mode, tex4ht has no robust support for minted: it relies on
% Pygments output that does not survive the htlatex DVI pipeline cleanly.
% We replace minted with verbatim-style fallbacks so the document still
% compiles. The PDF build keeps full minted (with -shell-escape).
% ============================
\ifSEhtml
  \usepackage{fvextra}
  % minted is NOT loaded in HTML mode. We provide a verbatim stand-in so
  % \begin{minted}[opts]{lang} ... \end{minted} still compiles. We use
  % fancyvrb's VerbatimEnvironment trick so the body is read verbatim
  % rather than tokenised by LaTeX (a +b body argument does NOT preserve
  % verbatim and breaks on '#', '%', '{', etc.).
  \newenvironment{minted}[2][]{%
    \VerbatimEnvironment
    \begin{Verbatim}[fontsize=\footnotesize]%
  }{%
    \end{Verbatim}%
  }
  % Inline replacement: \mintinline[opts]{lang}{code}
  % Note: \mintinline must accept verbatim-style content as its third arg.
  \NewDocumentCommand{\mintinline}{O{} m v}{\texttt{#3}}
  % Setting commands become no-ops in HTML mode
  \providecommand{\setminted}[2][]{}
  \providecommand{\setmintedinline}[2][]{}
\else
  \usepackage{minted}
\fi

\ifLuaTeX
  \usepackage{tagpdf}
  \tagpdfsetup{activate-all}
  \AtBeginEnvironment{tcolorbox}{\tagpdfparaOff}
  \AtEndEnvironment{tcolorbox}{\tagpdfparaOn}
\fi

% ============================
% Hyperref (load once, after most packages)
% ============================
\usepackage[
  colorlinks=true,
  linkcolor=blue,
  urlcolor=blue,
  citecolor=blue
]{hyperref}
\hypersetup{
  pdftitle={SE 423 Assignment},
  pdfauthor={Marius Juston},
}

% ============================
% unicode-math
% ----------------------------
% Only available on LuaTeX / XeTeX. NOT on pdfLaTeX. NOT on tex4ht.
% On tex4ht, attempting to load it triggers the luatexbase.sty error.
% ============================
\ifSEhtml\else
  \ifLuaTeX
    \usepackage{unicode-math}
  \else\ifXeTeX
    \usepackage{unicode-math}
  \fi\fi
\fi

\definecolor{highlight}{rgb}{0.396078431, 0, 0.721568627}

% pgfcalendar must be loaded in the preamble (here, in header.tex) rather
% than inside timline.tex. Loading \usepackage from within \input{...} is
% fragile under tex4ht (it triggers "Paragraph ended before \@reset@ptions").
\usepackage{pgfcalendar}

\input{timline}
\input{icons}

% ============================
% Assignment mode resolution
% ============================
\def\CourseCode{SE 423}
\def\CourseName{Mechatronics}

\def\AssignmentType{}
\def\AssignmentFooter{}

\makeatletter
\@namedef{Assignment@HW@type}{Homework Assignment}
\@namedef{Assignment@HW@footer}{HW}

\@namedef{Assignment@Lab@type}{Laboratory Assignment}
\@namedef{Assignment@Lab@footer}{Lab}

\@namedef{Assignment@LabVIEW@type}{LabVIEW Assignment}
\@namedef{Assignment@LabVIEW@footer}{LabVIEW}

\@ifundefined{Assignment@\AssignmentMode @type}{
  \PackageError{AssignmentMode}{Invalid AssignmentMode}{%
    Valid options: HW, Lab, LabVIEW}
}{
  \edef\AssignmentType{\@nameuse{Assignment@\AssignmentMode @type}}
  \edef\AssignmentFooter{\@nameuse{Assignment@\AssignmentMode @footer}}
}
\makeatother

% ============================
% Header / footer
% ============================
\pagestyle{fancy}
\fancyhf{}
\fancyfoot[L]{// \CourseCode, \CourseName}
\fancyfoot[R]{// \AssignmentFooter~\HWNum, Page \thepage\ of \pageref{LastPage}}
\renewcommand{\headrulewidth}{0pt}

% ============================
% Title box
% ============================
\newcommand{\MakeAssignmentTitle}{%
\par
\begin{tcolorbox}[
    colback=gray!25!white,
    width=\textwidth,
    boxrule=0.8pt,
    arc=2mm,
    left=6mm,
    right=6mm,
    top=4mm,
    bottom=4mm,
    center
]
\centering
\setstretch{1.2}

{\bfseries\Large \CourseCode\ \CourseName \\ \AssignmentType\ \#\HWNum\par}

\vspace{2mm}

{\bfseries
\AssignmentText\par
}
\end{tcolorbox}
\vspace{0.5cm}
}

% ============================
% minted styling (PDF mode only - stubs above are no-ops in HTML)
% ============================
\setmintedinline{breaklines, breakafter=_}

% Note: breaklines is intentionally OFF here. fvextra (which minted
% builds on) implements line wrapping via \everypar, and tagpdf's
% phase-III paragraph-hook counter desynchronises when a wrapped
% logical line crosses a page break — manifesting as a fatal
% "automatic begin/end text-unit para hooks differ" at \end{document}.
% Until tagpdf+fvextra cooperate cleanly under PDF/UA-2, leave wrapping
% off; long lines flow off the right margin rather than wrapping.
\setminted{
  breaklines=false,
  frame=lines,
  fontsize=\footnotesize,
}

\lstset{
    basicstyle=\ttfamily,   % Typewriter font for code
    breaklines=true,        % Allow breaking long lines
    columns=flexible        % Makes spacing better for inline
}

\newcounter{exercise}

\NewDocumentCommand{\Ex}{o}{%
  \par
  \stepcounter{exercise}
  \section*{Exercise \theexercise
    \IfValueT{#1}{: #1}
  }
  \addcontentsline{toc}{section}{Exercise \theexercise
    \IfValueT{#1}{: #1}}
}

\newcolumntype{C}[1]{>{\centering\arraybackslash}p{#1}}

\pgfplotsset{compat=1.18}

\newcommand{\MakeFooter}{%
\par
\begin{tcolorbox}[
    colback=gray!25!white,
    width=\textwidth,
    boxrule=0.8pt,
    arc=2mm,
    left=6mm,
    right=6mm,
    top=4mm,
    bottom=4mm,
    center
]
\centering
\setstretch{1.2}

{\bfseries Make sure that you \textcolor{highlight}{comment} you code and put your \textcolor{highlight}{initials} in your Gradescope submission! See \url{https://marius-juston.github.io/SE-423---Class-Material/Homeworks/commenting_code_example.pdf} for how we expect the code to be commented like!\par}
\end{tcolorbox}
}